\chapter{API remota, TUI e monitoraggio}\label{chap:remote}
Il servizio FastAPI in \pathcode{remote/} esegue via SSH tick idempotenti. Lo
stato autoritativo è \pathcode{.chain_state/} sul cluster; SQLite remoto è solo
cache/journal. Pause è soft e non cancella il job attivo. Tutte le route tranne
\pathcode{GET /} richiedono \pathcode{X-Auth-Token}.

\begin{lstlisting}[language=bash]
uv run --extra dev uvicorn app:app --port 8000 --app-dir remote
uv run --extra tui python remote/tui.py --url http://localhost:8000 --token local-token
\end{lstlisting}
Il registry espone celle canoniche e varianti manuali. Il campo storico API
\pathcode{ablation:true} seleziona/sostituisce la coda predefinita, non include
automaticamente PDA/hot. Per reward train il report qualificato è obbligatorio;
eval-only non deve riceverlo.
Rollout e Markov non sono queue config; il benchmark structured è assente da
driver e TUI. Un 401 indica autenticazione, mentre un 502 segnala tipicamente un
problema SSH/cluster; lo status può restituire cache marcata non raggiungibile.

\section{W\&B offline e sicurezza}
I compute scrivono run W\&B offline in directory isolate per gamba eval. Il sync
è successivo da ambiente con rete \cite{wandboffline}. Non mettere token HF,
W\&B, API o chiavi SSH in repository, YAML, log o command history condivisa;
preferire environment/secret store, permessi minimi e rotazione. Esporre l'API
solo su interfaccia fidata o tunnel SSH, usare token forte e non stampare header.
Validare ogni path come repository-relative e non accettare argomenti shell
arbitrari dalla TUI.

Da Windows gli script \pathcode{sync_cluster.ps1} scaricano log, checkpoint e
risultati. Conservare la gerarchia timestampata: appiattirla perde provenienza.
